\begin{chapintro}
架空の PCT の 10 研究を R に読み込み、研究ごとの感度と特異度を図で確かめます。カットオフ 0.5 ng/mL の 6 研究を取り出して、二項分布の二変量モデルを lme4 パッケージの glmer() で当てはめます。相関が境界に張り付いたら相関を 0 にしたモデルに単純化し、要約点と尤度比を出します。要約点と 2 つの領域の図の読み方を押さえ、全 10 研究には SROC 曲線だけを引きます。ROC 平面の描き方、十字プロット、reitsma() との比較、実データの AUDIT-C は、発展の印をつけて載せました。
\end{chapintro}

\section{データの準備と研究ごとの精度}
\label{sec:13-1}

\kkref{sec:12-10}で作った CSV ファイルを読み込みます。read.csv() や library() の使い方は\kkref{sec:05-5}で説明しました。列は研究名、場面（救急外来か ICU か）、カットオフ（ng/mL）、そして 2×2 表の 4 つの数 TP、FP、FN、TN です。mada パッケージは、列名が TP、FP、FN、TN になっていれば、そのまま 2×2 表として読みます。

\begin{rcode}
library(mada)    # 診断精度のメタ解析
library(lme4)    # 一般化線形混合モデル

pct <- read.csv("data/pct_bact.csv",
                fileEncoding = "UTF-8")
pct
\end{rcode}

最後の行で表を表示し、第 12 章の抽出結果と 1 行ずつ照合します。数を打ち間違えたまま解析しても、結果は正しくなりません。照合が済んだら、madad() で研究ごとの感度と特異度を計算します。correction.control = "single" は、2×2 表に 0 のセルがある研究だけ 0.5 を足して計算する設定です。この例には 0 のセルがないので、補正は働きません。全 10 研究は 3,626 人で、菌血症の人は 556 人です。

\begin{rcode}
d <- madad(pct, correction.control = "single")
print(d, digits = 2)
\end{rcode}

\begin{rout}
Diagnostic accuracies
      sens 2.5% 97.5% spec 2.5% 97.5%
 [1,] 0.62 0.48  0.75 0.80 0.76  0.84
 [2,] 0.56 0.44  0.68 0.78 0.74  0.82
 [3,] 0.79 0.64  0.89 0.75 0.67  0.81
  ...
[10,] 0.80 0.71  0.87 0.54 0.51  0.58
Test for equality of sensitivities:
X-squared = 30.514, df = 9, p-value = 0.00036
Test for equality of specificities:
X-squared = 154.9735, df = 9, p-value = <2e-16
\end{rout}

行番号はデータの並び順で、1 行目が Kato 2010、10 行目が Tanaka 2023 です。感度は 0.56 から 0.88、特異度は 0.54 から 0.91 まで散らばっています。その下の 2 つの検定は「すべての研究で感度（特異度）が等しい」という仮説の検定で、どちらも強く否定されています。この検定は出力に出ますが、判断には使いません。Cochrane DTA ハンドブックは、感度と特異度それぞれの等しさの検定を推奨していません（\kkref{sec:10-7}、Macaskill 2023）。研究間で精度がばらつくことは、次に描く図で見えます。

出力の最後には、感度と偽陽性割合の研究単位の相関（0.32）も表示されます。研究を点として並べただけの素朴な値なので、研究ごとの標本の大きさの違いを考慮していません。相関はあとで二変量モデルが推定し直します。

\subsection{対のフォレストプロット}
\label{sec:13-1-1}

感度と特異度を左右に並べたフォレストプロットを図 13.1 に示します。研究はカットオフの低い順に並べました。カットオフが 0.25 ng/mL の 2 研究（Nguyen 2014、Tanaka 2023）は感度が高く特異度が低く、カットオフが 1 ng/mL の Sato 2021 は特異度が 0.91 と最も高くなっています。介入のメタ解析と違い、この図の下に統合値の菱形を置くことはしません。感度と特異度を別々に統合すると、2 つの値の関係が失われるからです\kksee{sec:10-2}。

\kkfig[\linewidth]{p2_forest_pair.pdf}{図 13.1　研究ごとの感度と特異度（架空の PCT データ）。括弧内はカットオフ（ng/mL）、横線は 95\% 信頼区間}

mada の forest() に madad() の結果を渡すと、感度と特異度のフォレストプロットがそれぞれ描けます。

\begin{rcode}
forest(d, type = "sens")   # 感度のフォレストプロット
forest(d, type = "spec")   # 特異度のフォレストプロット
\end{rcode}

上の 2 行で描くと、感度と特異度は別々の図になります。図 13.1 は、この 2 つを左右に並べ、研究名とカットオフを左に添えて描いたものです。

横軸は感度（左）または特異度（右）で、0 から 1 の範囲です。四角が研究ごとの推定値、横線が 95\% 信頼区間です。人数の少ない研究ほど線が長くなります。左の感度の列では、病気のある人が 24 人の Park 2016 は 0.88（0.69〜0.96）と線が長く、95 人の Tanaka 2023 は 0.80（0.71〜0.87）と短くなっています。右の特異度の列では、 Tanaka 2023 の 0.54（0.51〜0.58）から Sato 2021 の 0.91（0.86〜0.94）まで開いていて、この 2 研究の区間は重なっていません。区間が重ならない研究の組があるのは、偶然のばらつきだけでは説明しにくい違いが研究の間にあることを示唆します。

\subsection{\hatten ROC 平面に研究を置く}
\label{sec:13-1-2}

次に、各研究を ROC 平面に 1 つの点として置きます（図 13.2）。横軸は偽陽性割合（1 − 特異度）、縦軸は感度です。点の大きさは研究の人数に比例させ、色はカットオフで分けました。左上に寄るほど精度の高い検査で、対角線上の点は当て推量と変わりません。

\kkfig[.68\linewidth]{p2_roc_points.pdf}{図 13.2　ROC 平面上の 10 研究。色が濃いほどカットオフが高い}

感度は TP ÷（TP + FN）、偽陽性割合は FP ÷（FP + TN）で計算できます。

\begin{rcode}
fpr_i <- pct$FP / (pct$FP + pct$TN)   # 偽陽性割合
sen_i <- pct$TP / (pct$TP + pct$FN)   # 感度
n <- with(pct, TP + FP + FN + TN)     # 研究の人数
pal <- c("#8FB0CE", "#2F5B82", "#14283C")
plot(NA, xlim = c(0, 1), ylim = c(0, 1),
     xlab = "1 − 特異度（偽陽性割合）", ylab = "感度")
abline(0, 1, lty = 3)                 # 当て推量の対角線
points(fpr_i, sen_i, pch = 21,
       cex = 0.8 + 2 * sqrt(n / max(n)),
       bg = pal[match(pct$cutoff, c(0.25, 0.5, 1))])
\end{rcode}

最後の bg の指定で、カットオフの値を 3 色のどれかに対応させています。

カットオフ 0.25 ng/mL の 2 研究（薄い点）は偽陽性割合が 0.29 と 0.46 で右寄りにあり、カットオフ 1 ng/mL の Sato 2021（濃い点）は偽陽性割合 0.09 で左端にあります。カットオフを下げると陽性と判定する人が増え、感度と偽陽性割合が一緒に上がるためで、これが閾値効果です\kksee{sec:10-1}。ただし点の散らばりはカットオフだけでは説明しきれません。同じ 1 ng/mL の Ono 2015 は偽陽性割合が 0.26 で 0.5 ng/mL の研究と同じあたりにあり、0.5 ng/mL の 6 研究だけでも感度は 0.56 から 0.88 まで開いています。

\subsection{\hatten 十字プロット}
\label{sec:13-1-3}

図 13.2 の点には、信頼区間がありません。研究ごとの不確かさを同じ平面に描き加えたものが、図 13.3 の十字プロットです。mada の crosshair() で描けます。

\kkfig[.68\linewidth]{x2_crosshair.pdf}{図 13.3　十字プロット。各研究の感度と偽陽性割合を点で示し、縦線が感度、横線が偽陽性割合の 95\% 信頼区間}

\begin{rcode}
crosshair(pct, xlim = c(0, 1), ylim = c(0, 1))
abline(0, 1, lty = 3)
\end{rcode}

軸は図 13.2 と同じで、十字の中心が 1 つの研究の推定値です。縦線の長さが感度の不確かさ、横線の長さが偽陽性割合の不確かさを表します。人数の少ない研究は十字が大きくなります。図 13.2 では点の大きさで人数を表しましたが、十字プロットでは線の長さで不確かさそのものを見せます。

線が表すのは研究の中の不確かさ（標本の大きさによる偶然のばらつき）で、研究どうしの違いではありません。研究間のばらつきは、十字の中心がどれだけ散らばっているかで見ます。この例では、10 研究のうち 9 研究の偽陽性割合が 0.09 から 0.29 の間に入る一方、偽陽性割合が約 0.46 の Tanaka 2023 は、短い線のまま右端に離れています。また、十字の線が重なっていても、研究どうしに違いがないとは言えません。感度と偽陽性割合の両方の線が重なる組は多くありますが、先に見た等しさの検定は、感度も特異度も研究間で等しくないことを示していました。

\section{カットオフ 0.5 ng/mL の 6 研究を取り出す}
\label{sec:13-2}

10 研究のカットオフは、0.25 ng/mL が 2 研究（Nguyen 2014、Tanaka 2023）、0.5 ng/mL が 6 研究（Kato 2010、Lopez 2011、Mori 2013、Park 2016、Qureshi 2018、Rossi 2019）、1 ng/mL が 2 研究（Ono 2015、Sato 2021）です。カットオフが混ざった研究から要約の感度と特異度を出すことは、どのカットオフの値でもなく臨床で解釈できないので避けます（\kkref{sec:10-6}、Macaskill 2023、9.3.1）。要約点は、研究がもっとも多いカットオフ 0.5 ng/mL の 6 研究だけから出します。

\begin{rcode}
pct5 <- subset(pct, cutoff == 0.5)
nrow(pct5)                                   # 6 研究
sum(with(pct5, TP + FP + FN + TN))           # 1782 人
sum(pct5$TP + pct5$FN)                       # 菌血症 272 人
\end{rcode}

元の pct は残したままで、全 10 研究は後で SROC 曲線とカットオフを共変量にした解析（\kkref{sec:14-2}）に使います。カットオフ 0.25 と 1 ng/mL は 2 研究ずつなので、それぞれの要約点は出しません。

\section{glmer() で二変量モデルを当てはめる}
\label{sec:13-3}

二変量モデルは、logit 感度と logit 特異度の組が研究ごとに 2 次元の正規分布からばらつく、と考えるモデルです。logit は\kkref{sec:10-3}、二変量モデルの考え方は\kkref{sec:10-4}、研究ごとのずれをランダム効果で表す考え方は\kkref{sec:05-4}で説明しました。研究内のばらつきは二項分布のまま扱います（Chu 2006）。この方法は 0 のセルがある研究でも補正なしに扱え、Cochrane の DTA ハンドブック第 10 章も、R では lme4 パッケージの glmer() で当てはめる例を示しています。glmer() は一般化線形混合モデル（研究内は二項分布、研究間は正規分布という 2 段のモデル）を当てはめる関数です。

感度は病気のある人から、特異度は病気のない人から計算する別々のデータです。1 つの式で両方を扱うために、1 研究を「感度の行」と「特異度の行」の 2 行に分けた縦長のデータを作ります。to\_long() がこの縦長のデータを作る関数です。sens の行では y が真陽性の数、n が疾患ありの人数、spec の行では y が真陰性の数、n が疾患なしの人数になります。

\begin{rcode}
to_long <- function(d) rbind(
  data.frame(study = d$study, sens = 1, spec = 0,
             y = d$TP, n = d$TP + d$FN),
  data.frame(study = d$study, sens = 0, spec = 1,
             y = d$TN, n = d$TN + d$FP))
long5 <- to_long(pct5)
g2 <- glmer(cbind(y, n - y) ~ 0 + sens + spec +
              (0 + sens + spec | study),
            data = long5, family = binomial)
summary(g2)
\end{rcode}

式を部品に分けると、表 13.1 のとおりです。初めは式を写して実行できれば十分で、意味は結果を読んだあとで戻って確かめれば足ります。

\begin{kkcenter}
\kkfit{\footnotesize\begin{tabular}{p{40mm}p{66mm}}
\toprule
部品 & 意味 \\
\midrule
cbind(y, n − y) & 正しく判定された人数 y と、誤って判定された人数 n − y の組。sens の行は真陽性と偽陰性、spec の行は真陰性と偽陽性 \\
\textasciitilde{} & 左が説明される結果、右がそれを説明する項 \\
0 + sens + spec & 切片を置かず、logit 感度と logit 特異度を別々の平均として推定する \\
(0 + sens + spec | study) & 研究ごとのずれを、感度と特異度の両方について推定する。2 つのずれの相関も推定する \\
family = binomial & 人数を二項分布で扱う \\
\bottomrule
\end{tabular}}
\kkcap{表 13.1　glmer() の式の部品}
\end{kkcenter}

実行すると、glmer() の時点で boundary (singular) fit という警告が出て、summary() の末尾にも同じ表示が出ます。エラーではなく、推定が止まったわけでもありません。相関や標準偏差が、取りうる範囲の端（相関なら −1 か 1、標準偏差なら 0）に張り付いた、という知らせです。

出力で見る場所は 2 つです。Random effects の Std.Dev. と Corr は、logit の尺度での研究間の標準偏差と、感度と特異度のずれの相関です。Variance は Std.Dev. を 2 乗した値で、同じ情報の別の表し方です。Fixed effects の Estimate は、logit 感度（sens）と logit 特異度（spec）の平均です。

\begin{rout}
Random effects:
 Groups Name Variance Std.Dev. Corr
 study  sens 0.25253  0.5025
        spec 0.01978  0.1407   -1.00
Fixed effects:
     Estimate Std. Error z value
sens  1.06422    0.25174   4.228
spec  1.11447    0.08379  13.301
boundary (singular) fit: see help('isSingular')
\end{rout}

Corr の欄で、感度と特異度の相関が −1.00 になっています。相関は −1 から 1 までしかとれないので、これ以上は動けない端に張り付いた状態です。研究が 6 本と少なく、2 つのばらつきと相関を同時に決めるだけの情報が足りないときに、このような推定になります。端に張り付いた推定では、標準誤差と信頼区間の計算が不安定になり、要約点の値も不安定になります。この結果から出した感度 0.743 と特異度 0.753 は、参考にとどめ、報告には使いません。

\section{相関を 0 にしたモデルに単純化する}
\label{sec:13-4}

Cochrane の DTA ハンドブックは、二変量モデルが収束しない、相関が ±1 に張り付く、分散が 0 になるといったときは、モデルを単純化するよう勧めています（Macaskill 2023、9.4.8、Takwoingi 2023、10.6.2）。単純化するかどうかは、研究の本数ではなく推定の様子で決めます。典型的な単純化が、感度と特異度の相関を 0 にするモデルです（相関の意味は\kkref{sec:10-4}で説明しました）。式の上では、2 つのランダム効果を別々の括弧に分けます。

\begin{rcode}
g <- glmer(cbind(y, n - y) ~ 0 + sens + spec +
             (0 + sens | study) + (0 + spec | study),
           data = long5, family = binomial)
summary(g)
\end{rcode}

「(0 + sens | study) + (0 + spec | study)」は、感度と特異度が研究ごとに別々にばらつき、その間の相関は考えないという式です。今度は singular の警告は出ません。

\begin{rout}
Random effects:
 Groups  Name Variance Std.Dev.
 study   sens 0.26801  0.5177
 study.1 spec 0.01841  0.1357
Fixed effects:
     Estimate Std. Error z value
sens   1.0773     0.2615    4.12
spec   1.1353     0.0832   13.64
\end{rout}

研究間の標準偏差は、logit 感度が 0.518、logit 特異度が 0.136 です。感度は研究によって大きくばらつき、特異度はあまりばらつきません。6 研究では、相関を推定するだけの情報がありません。相関を 0 と置くのは、感度と特異度の研究ごとのずれを別々に推定するという簡単化で、要約点への影響は小さいことが多いとされます（ハンドブック 10.6.2、Takwoingi 2023）。代わりに、相関を使って精度を高める利点は失い、感度と特異度が一緒に動く関係は要約点の不確かさに入りません。この制約は報告に書きます。以降の要約点は、このモデルのものを使います。

\begin{hattenbox}
glmer() は、ランダム効果が 2 つ以上あると、数値積分の点（求積点）を増やせず、ラプラス近似で推定します（Macaskill 2023、10.2.3）。研究が少ないときの推定が不安定になりやすい一因です。
\end{hattenbox}

\begin{KKfbox9ao}[R を使わない場合]
R を使わずに二変量モデルを当てはめたいときは、無料の Web ツール MetaDTA が使えます。MetaDTA は内部で glmer() を使っていて、2×2 表の CSV ファイルを読み込ませれば要約点と SROC 曲線が出ます。
\end{KKfbox9ao}

\section{要約点と尤度比}
\label{sec:13-5}

推定値を逆 logit で戻して、要約点とその 95\% 信頼区間を出します。信頼区間は logit の尺度で正規近似し、戻します。

\begin{rcode}
b <- fixef(g)
V <- as.matrix(vcov(g))
se_b <- sqrt(diag(V))
plogis(cbind(est = b, lower = b - 1.96 * se_b,
             upper = b + 1.96 * se_b))
\end{rcode}

\begin{rout}
           est     lower     upper
sens 0.7459900 0.6375786 0.8305890
spec 0.7568185 0.7255653 0.7855644
\end{rout}

カットオフ 0.5 ng/mL の PCT の、要約の感度は 0.746（95\% 信頼区間 0.638〜0.831）、特異度は 0.757（0.726〜0.786）です。特異度の区間が狭いのは、特異度のばらつきが小さく、疾患のない人の数も多いからです。

\begin{KKfbox9ao}[R のメモ]
logit は、割合 p を log(p / (1 − p)) に変える変換です。0 から 1 の値を、負の無限大から正の無限大までの値に広げます。R では qlogis() で logit に、plogis() で元の割合に戻せます。たとえば plogis(1.0773) は 0.746 になり、上の出力の感度と一致します。詳しくは\kkref{sec:10-3}を読み直してください。
\end{KKfbox9ao}

要約点から、陽性尤度比、陰性尤度比、診断オッズ比を計算します\kksee{sec:09-4}。

\begin{rcode}
se <- unname(plogis(b["sens"]))
sp <- unname(plogis(b["spec"]))
c(LRpos = se / (1 - sp), LRneg = (1 - se) / sp,
  DOR = (se / (1 - se)) / ((1 - sp) / sp))
\end{rcode}

\begin{rout}
    LRpos     LRneg       DOR
3.0676260 0.3356287 9.1399395
\end{rout}

陽性尤度比は 3.07、陰性尤度比は 0.34、診断オッズ比は 9.1 でした。陽性尤度比 3.07 は、PCT が陽性の人が菌血症である見込み（オッズ）が、検査前の 3.07 倍になるという意味です。陰性尤度比 0.34 は、陰性なら見込みが 3 分の 1 ほどになるという意味です。どちらも、検査の結果だけで菌血症を確定したり否定したりできる大きさではありません（\kkref{sec:15-1}で 1000 人あたりの数にします）。

尤度比を研究ごとに計算して別々に統合する方法は勧められていません。統合した陽性尤度比と陰性尤度比の組が、ありえない感度と特異度の組に対応してしまうことがあるからです（Zwinderman 2008、Macaskill 2023、9.3.1）。要約点から計算する、という順序を守ります。要約の尤度比の信頼区間は、この冊子では計算しません。

\section{要約点と 2 つの領域を図にする}
\label{sec:13-6}

要約点と、その不確かさを表す領域を図 13.4 に示します。赤い菱形が要約点（感度 0.746、偽陽性割合 0.243）、濃い楕円が 95\% 信頼領域、淡い大きな楕円が 95\% 予測領域、丸がカットオフ 0.5 ng/mL の 6 研究です。

\kkfig[.7\linewidth]{p2_summary05.pdf}{図 13.4　カットオフ 0.5 ng/mL の 6 研究の要約点、95\% 信頼領域、95\% 予測領域}

信頼領域は、要約点（感度と特異度の平均）がどこにあるかの不確かさを表します。研究を増やせば小さくなります。予測領域は、これから新しく行う研究の感度と特異度がどのあたりに来そうかを表し、研究間のばらつきそのものを含むので、研究を増やしても小さくなりません。介入のメタ解析の信頼区間と予測区間の関係（\kkref{sec:06-3}）を 2 次元にしたものです。

図 13.4 の信頼領域は小さく、感度の方向にやや長い形です。予測領域は縦に細長く、感度の方向に広く、偽陽性割合の方向には狭くなっています。感度の研究間標準偏差が 0.518、特異度が 0.136 と、感度のほうがはるかにばらつくためです。ある病院で PCT をカットオフ 0.5 ng/mL で導入したとき、その病院の特異度は要約点の近くに収まりそうですが、感度は要約点から大きく外れてもおかしくない、という意味です。要約点だけを報告すると、この広がりが読み手に伝わりません。

予測領域は、研究間の分布が 2 次元の正規分布に従うという仮定に頼っています。研究数が少ないと、領域は実際より広くも狭くも出ます（\kkref{sec:10-7}、Macaskill 2023、9.3.2）。6 研究の図 13.4 は、形と大きさを目安として読みます。

\subsection{\hatten 図 13.4 を描くコード}
\label{sec:13-6-1}

楕円は、logit の尺度で要約点のまわりに作り、ROC 平面に戻して描きます。信頼領域は要約点の推定値の分散から、予測領域は研究間のばらつきを足した分散から作ります。ellipse\_roc() は、楕円の点を計算する関数です。中身は読まなくてかまいません。

\begin{rcode}
ellipse_roc <- function(mu, S, n = 200) {
  r <- sqrt(qchisq(0.95, 2))
  th <- seq(0, 2 * pi, length.out = n)
  e <- t(mu + r * t(chol(S)) %*% rbind(cos(th), sin(th)))
  data.frame(fpr = plogis(e[, 2]), sens = plogis(e[, 1]))
}
\end{rcode}

この関数を使って、要約点と 2 つの領域を描きます。

\begin{rcode}
vc <- VarCorr(g)
S_btw <- diag(c(attr(vc[[1]], "stddev")^2,
                attr(vc[[2]], "stddev")^2))
mu <- c(b["sens"], -b["spec"])  # logit 感度、logit 偽陽性割合
V_m <- matrix(c(V[1, 1], -V[1, 2], -V[1, 2], V[2, 2]), 2)
conf <- ellipse_roc(mu, V_m)           # 信頼領域
pred <- ellipse_roc(mu, S_btw + V_m)   # 予測領域
plot(NA, xlim = c(0, 1), ylim = c(0, 1),
     xlab = "1 − 特異度（偽陽性割合）", ylab = "感度")
abline(0, 1, lty = 3)
polygon(pred$fpr, pred$sens, border = "grey50", lty = 2)
polygon(conf$fpr, conf$sens)
points(fpr(pct5), sens(pct5), pch = 21)
points(1 - sp, se, pch = 23, bg = "red")
\end{rcode}

\section{\hatten mada の reitsma() との比較}
\label{sec:13-7}

同じ 6 研究を mada の reitsma() で解析することもできます。研究内を正規近似する方法で、参考として示します。

\begin{rcode}
summary(reitsma(pct5))
\end{rcode}

\begin{rout}
sensitivity          0.737  0.624  0.826
false pos. rate      0.248  0.217  0.282
Variance components:
      Std. Dev tsens  tfpr
tsens    0.543 1.000     .
tfpr     0.159 1.000 1.000
\end{rout}

要約の感度は 0.737（0.624〜0.826）、偽陽性割合は 0.248（0.217〜0.282）で、glmer() の結果とほとんど変わりません。一方、研究間の相関の行は 1.000 で、境界に張り付いたことは reitsma() でも同じです（偽陽性割合で表すので、符号は glmer() と逆になります）。reitsma() は相関を 0 に固定する指定ができないので、単純化したモデルは glmer() で当てはめます。

Cochrane の DTA ハンドブック第 10 章は、研究内のばらつきに二項分布を使うことを勧め、正規近似の mada は Cochrane レビューには推奨しないとしています（Takwoingi 2023、10.2.3）。理由として、疾患ありと疾患なしの人数が十分に多いという近似の前提が崩れると推定が偏ることと、0 のセルに補正を足す必要があることを挙げています。二変量モデルの 2 つの階と、研究内の扱い方の違いは発展の\kkref{sec:10-4-1}で説明しました。本書でも glmer() を主解析とし、reitsma() は SROC 曲線を引くときや、結果を確かめるときに使います。

\section{全 10 研究の SROC 曲線}
\label{sec:13-8}

カットオフが混ざった全 10 研究には、SROC 曲線だけを引きます\kksee{sec:10-6}。mada の reitsma() は、研究内を正規近似する二変量モデルを当てはめる関数です。reitsma() で全 10 研究にモデルを当てはめ、sroc() で曲線の座標を取り出して描きます。

\begin{rcode}
fit_all <- reitsma(pct)
plot(NA, xlim = c(0, 1), ylim = c(0, 1),
     xlab = "1 − 特異度（偽陽性割合）", ylab = "感度")
abline(0, 1, lty = 3)
lines(sroc(fit_all), lwd = 2)       # 曲線だけを描く
points(fpr(pct), sens(pct), pch = 21)
\end{rcode}

summary(fit\_all) の出力には要約点（感度と偽陽性割合）も表示されますが、カットオフが混ざっているので使いません。曲線だけを使います。

\kkfig[.7\linewidth]{p2_sroc.pdf}{図 13.5　全 10 研究の SROC 曲線。要約点は描かない。点の色はカットオフ}

図 13.5 の曲線は、カットオフを動かしたときに感度と偽陽性割合がどう入れ替わるかの平均的な関係を表します。点は研究ごとの感度と偽陽性割合で、曲線のまわりに散らばっています。

曲線は左下から右上まで引かれていますが、データがあるのは偽陽性割合がおよそ 0.09 から 0.46 の範囲だけです。その外側の曲線は、モデルの形をそのまま延ばした線にすぎません。たとえば「偽陽性割合 0.05 なら感度はいくつか」を曲線から読むと、データのない場所の値を読むことになります。曲線は観測された点の範囲の中だけで読みます。

summary(fit\_all) の末尾には AUC（0.824）も表示されます。AUC は SROC 曲線の下の面積で、参考値として扱います。観測された点の範囲の外に延ばした曲線の部分も面積に入るため、データが支えていない領域の値を含み、特定のカットオフでの感度と特異度も表さないからです。mada は観測された偽陽性割合の範囲に限った部分 AUC（0.755）も出します。報告には、曲線そのものと、要約点を出せる場合はその点を示します。

\section{\hatten 実データの AUDIT-C で再現する}
\label{sec:13-9}

mada パッケージには、AUDIT-C（飲酒の問題を見つけるための 3 問の質問票）の 14 研究のデータが AuditC として入っています（Kriston 2008）。このデータは研究ごとの 2×2 表だけを持ち、カットオフの情報は含みません。二変量モデルの手順を実データでたどる例として、14 研究をそのまま統合します。

\begin{rcode}
data("AuditC")
ac <- as.data.frame(AuditC)
ac$study <- seq_len(nrow(ac))
ga <- glmer(cbind(y, n - y) ~ 0 + sens + spec +
              (0 + sens + spec | study),
            data = to_long(ac), family = binomial)
summary(ga)
plogis(fixef(ga))
\end{rcode}

14 研究では、相関が境界に張り付くことなく推定でき、singular の警告も出ません。

\begin{rout}
Random effects:
 Groups Name Variance Std.Dev. Corr
 study  sens 2.5959   1.6112
        spec 0.3812   0.6174   -0.82
\end{rout}

感度と特異度の相関は −0.82 で、閾値効果の向きに強く負です。要約の感度は 0.925（95\% 信頼区間 0.834〜0.968）、特異度は 0.782（0.720〜0.833）でした。感度と特異度を別々に統合すると感度 0.930、特異度 0.783 になり、点推定はほとんど変わりません。二変量モデルの利点は、この相関を推定して信頼領域や SROC 曲線に反映できる点にあります。

同じデータを reitsma() で解析すると、感度は 0.891（0.808〜0.941）、偽陽性割合は 0.220（0.167〜0.285）、AUC は 0.887 でした。図 13.6 に SROC 曲線と信頼領域を示します。

\kkfig[.65\linewidth]{p2_auditc_sroc.pdf}{図 13.6　AUDIT-C の 14 研究の SROC 曲線と 95\% 信頼領域}

\begin{rcode}
fa <- reitsma(AuditC)
plot(fa, extrapolate = TRUE)
points(fpr(AuditC), sens(AuditC))   # 14 研究の点
\end{rcode}

図の読み方は PCT の図 13.4 と同じです。要約点が感度 0.89、偽陽性割合 0.22 の位置にあり、研究の点は感度の方向に広く散らばっています。95\% 信頼領域は小さく見えますが、予測領域を描いていないので、新しい研究の値がどこに来るかは、この図からは分かりません。

このデータには感度が 1.00 に近い研究があり、2×2 表に 0 のセルが 2 つあります。reitsma() は 0 のセルに 0.5 を足して計算しました。glmer() の感度は reitsma() より 0.03 ほど高くなっています。0.5 を足す補正が、感度が 1 に近い研究の値を引き下げたためと考えられます。PCT の 6 研究で両者が近く、AUDIT-C で食い違ったのは、0 のセルがあり感度が 1 に近いという、近似の前提が崩れる条件に AUDIT-C が当たるためと考えられます。

\section{この章のまとめ}
\label{sec:13-10}

検査のメタ解析では、研究ごとの感度と特異度を表と図で確かめてから、要約点を出す研究の範囲を決めます。例題では、カットオフが混ざった全 10 研究から要約点を出さず、共通のカットオフ 0.5 ng/mL の 6 研究から二項分布の二変量モデルで要約点を推定しました。glmer() の式は、写して実行できれば初めは十分です。見る場所は、Random effects の Std.Dev. と Corr、Fixed effects の Estimate です。相関が −1 の境界に張り付いたので、相関を 0 にしたモデルに単純化し、要約の感度 0.746、特異度 0.757、陽性尤度比 3.07、陰性尤度比 0.34 を得ました。全 10 研究には SROC 曲線だけを引き、観測範囲の中で読みます。予測領域は感度の方向に広く、研究数が少ないので目安にとどめます。\kkchref{14}では、全 10 研究にカットオフを共変量として入れ、ばらつきの原因を探します。
